\textbf{Línea base de capacidad y escalado.}
La Tabla~\ref{tab:capacidad-cloud} fija recursos por tarea y cantidades por unidad desplegable. Normal y peak corresponden a 25 y 100 HTTP/s; crecimiento corresponde a 300 HTTP/s y tres veces los eventos de la mezcla de SD5, 5.4.1. Son configuraciones seleccionadas de diseño, sujetas a aceptación con carga representativa. La memoria incluye aplicación, agente de protección y observabilidad: no se añaden esos agentes fuera del presupuesto de la tarea. Las combinaciones de 1 vCPU/4 GiB y 2 vCPU/8 GiB son admitidas por Fargate; esa compatibilidad no acredita rendimiento de la aplicación \parencite{c8-fargate-recursos}.

\begin{table}[htbp]
\centering
\fontsize{9}{11}\selectfont
\begin{tabularx}{\linewidth}{|X|r|r|r|r|}
\hline
\EncabezadoTabla{Unidad} & \EncabezadoTabla{vCPU/GiB} & \EncabezadoTabla{Normal} & \EncabezadoTabla{Peak} & \EncabezadoTabla{Crec./máx.}\\ \hline
Marina-Ops / Marina-Log / Marina-Admin (cada uno) & 1/4 & 2 & 4 & 6/8\\ \hline
Integración externa & 1/4 & 2 & 4 & 6/8\\ \hline
Notificaciones & 1/4 & 2 & 4 & 4/8\\ \hline
Telemetría & 1/4 & 2 & 2 & 6/8\\ \hline
Sincronización & 2/8 & 2 & 4 & 6/8\\ \hline
Auditoría & 1/4 & 2 & 2 & 6/8\\ \hline
\end{tabularx}
\caption{Tareas seleccionadas por unidad y escenario cloud}
\label{tab:capacidad-cloud}
\FuenteTabla{Cálculo de Synaptix; SD5, 5.4.1; BT RT-02.10 y RT-15.01. Cantidades por unidad, no por zona.}
\end{table}

La base normal suma $3\times2+5\times2=16$ tareas, 18 vCPU y 72 GiB; peak suma 28 tareas, 32 vCPU y 128 GiB; crecimiento suma 46 tareas, 52 vCPU y 208 GiB. El techo de ocho tareas para cada una de las ocho unidades suma 64 tareas, 72 vCPU y 288 GiB. Se reparten por igual entre dos zonas. El techo incluye el estado estable; una liberación permite hasta el doble transitorio mediante \texttt{maximumPercent=200}, con capacidad y cuotas reservadas, y devuelve luego la cantidad al rango autorizado. El inventario y la valorización económica distinguen consumo estable, liberaciones y ensayos.

La capacidad de las tres unidades conserva una hipótesis de reparto agregado por solicitudes, no por número de módulos. La prueba incorpora combustible y consumos en Marina-Ops y residuos en Marina-Log. Para los núcleos se adopta una mezcla de solicitudes del 50/30/20\,\% entre Marina-Ops, Marina-Log y Marina-Admin, y una hipótesis de 0,010 segundos de CPU por solicitud. Se reserva por tarea el 60\,\% de CPU para negocio, el 10\,\% para agentes y el 30\,\% libre. En el núcleo más cargado, peak demanda $100\times0,50\times0,010=0,50$ vCPU y crecimiento $300\times0,50\times0,010=1,50$. Tres tareas supervivientes de las seis de crecimiento aportan $3\times0,60=1,80$ vCPU para negocio. La memoria se distribuye inicialmente en 2 GiB de heap, 1 GiB de memoria nativa/agentes y 1 GiB de reserva por tarea de núcleo. El ensayo debe comprobar esas hipótesis y las latencias P95 de SD5; si no se cumplen, se aumenta recurso o cantidad antes de habilitar la carga. La fórmula no convierte un supuesto de CPU en benchmark.

El núcleo Marina-Admin incluye los trabajos propios MASS-5 y ENV-5 de SD5, además del controlador HTTP de ANAL-5. Se selecciona un cupo global de ejecutor pesado base y tres en crecimiento, con lease durable y caudal compartido entre MASS/ENV/EXP y extracciones ANAL cuando comparten tarea; no se añade ese cupo por usuario, destinatario o réplica. Con hipótesis de validación de diez filas/s a 0,020 s CPU/fila, aplicación de dos registros/s a 0,010 s CPU/registro y serialización de 0,020 vCPU, el cupo base utiliza 0,24 vCPU; entrega/descarga compartida de 2 MiB/s a 0,020 s CPU/MiB añade 0,04. Marina-Admin demanda 0,33 vCPU en normal y 0,48 en peak con una tarea superviviente, frente a 0,60 para negocio; crecimiento requiere $0,60+3(0,24+0,04)=1,44$ vCPU, frente a 1,80 en tres supervivientes. Buffers de 128 MiB y temporal de 256 MiB se comparten en el presupuesto existente, sin reservas independientes por módulo. Los 170/510 GiB de objetos MASS/ENV son una reserva adicional separada de datos y journal: 10/30 para cargas y 160/480 para informes/copias nominativas. Se respetan cortes, permisos, retención de envíos pendientes y mínimos operacionales; la cola limita analítica/intercambio antes de comprometer captura. Todas las hipótesis se validan con carga, reconexión y una zona retirada; una cola que no termina o una demanda mayor exige ampliar antes de habilitar. El dimensionamiento integral y la valorización OE siguen pendientes; estos subtotales no los acreditan.

Para integración, telemetría y auditoría se adopta 0,005 segundos de CPU por mensaje lógico procesado. Una tarea superviviente admite como objetivo $0,60/0,005=120$ mensajes/s; el crecimiento de integración externa, $111\times3=333$ mensajes/s, exige al menos tres supervivientes. Integración prueba conjuntamente los peaks de sus fuentes, y telemetría prueba también las ráfagas locales y centrales de SD5. Notificaciones tiene ejecutores separados por canal y una reserva de cinco envíos/s por canal, con el volumen y los reintentos del catálogo de integración. Sincronización fija como condición de aceptación cuarenta MB/s equivalentes por tarea superviviente: una satisface el lote base y tres el objetivo de 120 MB/s de crecimiento. Se verifica la cadena completa de LAN, archivos, red, validación y persistencia; sumar capacidades de tareas no acredita que un cuello de botella local permita alcanzarlas.

El escalado automático de núcleos sigue objetivos de CPU de aplicación del 60\,\% y memoria total del 70\,\%. La métrica de CPU de aplicación excluye agentes, se publica mediante OpenTelemetry y alimenta la política; se vigila además CPU total, cuyo presupuesto incluye el 10\,\% de agentes. Los workers agregan tareas si la edad del mensaje más antiguo supera treinta segundos o hay más de cien mensajes pendientes por tarea durante dos minutos; el vaciado de emergencia activa cuatro tareas de notificación antes del envío. Los núcleos también aumentan ante P95 por encima del umbral del tipo de operación durante dos ventanas de un minuto. El enfriamiento de ampliación es de sesenta segundos y el de reducción de quince minutos. No se reduce mientras haya reconexión, emergencia, despliegue o pérdida de zona; nunca se cruza el mínimo de dos ni el máximo de ocho. Antes de temporada alta se programa peak; antes de habilitar crecimiento se establece su cantidad prevista y se ejecuta la prueba de pérdida de zona. Un máximo alcanzado con latencia o cola creciente genera incidencia de capacidad: se amplía la línea base mediante control de cambios, sin rechazar hechos ya aceptados \parencite{c8-ecs-autoescalado}.

\textbf{Ambientes y recuperación.}
Desarrollo y QA utilizan una tarea de cada unidad, nueve vCPU y 36 GiB por ambiente, con datos sintéticos y separación de permisos. Se planifican 160 y 120 horas mensuales de uso, respectivamente, como supuestos de consumo y no como límites de trabajo contractual. Preproducción usa la topología de dieciséis tareas durante las ventanas de ensayo, inicialmente ochenta horas mensuales, y escala al escenario que se comprueba. Sus ensayos de falla requieren ambas zonas; no se extrapola el resultado de una instancia de QA. Producción y DR conservan funcionamiento 24\,$\times$\,7. DR mantiene las dieciséis tareas mínimas distribuidas en dos zonas, sin escritura de negocio hasta la promoción, y recibe la cantidad de peak o crecimiento que rija en Producción antes de cambiar tráfico. La capacidad se prueba en DR; no se presupone por replicar imágenes. Las cuotas regionales reservan también el doble transitorio del techo para una liberación.

Se selecciona RDS PostgreSQL 18 \texttt{db.m7g.large}, 2 vCPU/8 GiB, para cada uno de los dos almacenes primarios y sus réplicas regionales, con Multi-AZ. El transaccional asigna 200 GiB físicos y el de proyecciones 50 GiB; las reservas lógicas de 200 GB y 15 GB de SD5 no son unidades de provisión del proveedor. Se elige gp3: en esos tamaños AWS documenta 3.000 IOPS y 125 MiB/s de referencia del volumen. El transaccional amplía hasta 1.024 GiB y las proyecciones hasta 200 GiB. Se registra espacio libre con alertas al 30 y 20\,\%, evitando alcanzar el umbral de ampliación sin capacidad de seguimiento. El crecimiento selecciona \texttt{db.m7g.2xlarge}, 8 vCPU/32 GiB, y al menos 400 GiB gp3 para cada almacén: la configuración del volumen ofrece 12.000 IOPS y 500 MiB/s de referencia, sujetos al límite de la instancia y al comportamiento real. La réplica DR se amplía antes que su primario. Clase, motor y disponibilidad conjunta se comprueban en ambas regiones antes de provisión; un catálogo de clase o de disco no acredita los 40/120 MB/s del proceso \parencites{c8-rds-clases,c8-rds-almacenamiento}.

\textbf{Equivalencia de Preproducción.}
La Tabla~\ref{tab:preprod-equivalencia} fija la comparación del perfil base y del crecimiento habilitado. SRE mantiene una exportación normalizada de IaC, políticas, contratos, parámetros de base de datos y manifiestos por ambiente; Calidad compara versión, digest y configuración antes de cada ensayo y promoción. Se excluyen de igualdad únicamente identificadores de cuenta, región cuando corresponda, nombres, direcciones, claves/secretos propios y destinos de prueba. Cada diferencia queda identificada, justificada y aprobada en \texttt{EQ-4}, versión 1.0; una diferencia adicional sin aprobación bloquea promoción. Los valores de seguridad, autorización, tiempos de espera, reintentos, límites de negocio y retención mantienen el perfil productivo.

\begingroup\setlength{\tabcolsep}{3pt}
\begin{longtable}{L{\dimexpr.20\textwidth-2\tabcolsep\relax}L{\dimexpr.38\textwidth-2\tabcolsep\relax}L{\dimexpr.42\textwidth-2\tabcolsep\relax}}
\caption{Producción y Preproducción: igualdad y diferencias seleccionadas}\label{tab:preprod-equivalencia}\\\hline
\EncabezadoTabla{Componente} & \EncabezadoTabla{Producción} & \EncabezadoTabla{Preproducción y justificación}\\\hline\endfirsthead
\hline\EncabezadoTabla{Componente} & \EncabezadoTabla{Producción} & \EncabezadoTabla{Preproducción y justificación}\\\hline\endhead
ECS/Fargate & Ocho unidades; dos zonas; 16/28/46 tareas en normal/peak/crecimiento; techo estable 64. & Mismas ocho unidades, recursos por tarea, zonas y políticas. 16 tareas en validación ordinaria; adopta 28 o 46 al ensayar el escenario respectivo y conserva techo 64 y doble transitorio de liberación.\\\hline
RDS PostgreSQL & Dos almacenes Multi-AZ; PostgreSQL 18; base 2 vCPU/8 GiB, 200/50 GiB gp3; crecimiento 8 vCPU/32 GiB y al menos 400 GiB por almacén. & Mismos motor/minor del manifiesto, clase, parámetros, extensiones, almacenamiento e IOPS/caudal del escenario ensayado. No se usan datos ni claves productivos.\\\hline
Red y exposición & Dos zonas, rutas redundantes, VPC Link, API Gateway, balanceo, WAF, identidad y permisos separados. & Misma topología y políticas; cuentas, DNS, certificados, direcciones y principales propios. Tráfico y permisos aislados para impedir modificar operación o destinos reales.\\\hline
Eventos, objetos y observabilidad & Mismos contratos de colas/eventos, DLQ, idempotencia, versionado/retención S3 y controles de registros. & Igual configuración funcional y de seguridad; datos sintéticos representativos, IDs y buckets propios. Igual volumen y distribución del ensayo, incluida retención reconstruida para medir consultas.\\\hline
Integraciones externas & Adaptadores, contratos, cuotas contratadas y límites del catálogo A4.2. & Mismo código; endpoints de prueba y emuladores controlados evitan operaciones/avisos reales. Se reproducen tiempos, fallos, cuotas y ráfagas; se verifica aparte contra el endpoint real autorizado antes de habilitar. Un emulador no acredita capacidad o contrato del proveedor.\\\hline
Célula local & Dos nodos R360 de T-11, ítem 32, réplica síncrona y único escritor; par de cortafuegos y core redundante. & Laboratorio segregado: dos R360, dos FG-120G y dos C9300-24P-E con la configuración de los ítems 32/34/36; mismo sistema, versiones, almacenamiento, HA, rutas A/B y políticas. T-11, ítem 32P, separa estos seis equipos del parque y repuestos productivos.\\\hline
Disponibilidad y datos & Producción 24\,$\times$\,7; datos reales y capacidad habilitada. & Ventanas ordinarias de 80 h/mes como supuesto de consumo; se amplían para cada ensayo. Datos protegidos con la misma cardinalidad, tamaños, distribución y restricciones del escenario. Ahorra horas de uso; no reduce recursos durante aceptación ni sustituye pruebas físicas del recinto real.\\\hline
\end{longtable}
\FuenteTabla{Decisiones de Synaptix: línea base de esta sección, T-11 y SD5, 5.4.1; BT RT-04.02, p.~10.}
\endgroup

Antes de cada paso a Producción se reserva capacidad y cuotas para carga y estrés a $1,5$ veces el peak declarado: en el perfil inicial, $1,5\times100=150$ HTTP/s, además de la mezcla concurrente de transacciones, eventos, reconexión y avisos de SD5/A4.2. Se ensaya también crecimiento a 300 HTTP/s; si éste pasa a ser el peak habilitado, el ensayo correspondiente exige 450 HTTP/s y actualización previa del perfil donde resulte necesaria. No se extrapola el resultado de 150 HTTP/s a ese peak nuevo. SRE provisiona el perfil elegido y verifica recursos reales antes de ejecutar; Calidad mide latencias, errores, integridad, colas, pérdida de tarea/zona y retorno. Si no supera los umbrales, se ajusta la misma línea base de Producción y Preproducción y se repite; no se aumenta sólo el laboratorio para ocultar un límite productivo. La matriz EQ-4, los manifiestos y el informe efectivo se incorporan a LIB-6. Las funciones locales se ensayan en el par segregado con carga representativa y desconexión; la aceptación de energía y equipos de campo se verifica además en su instalación real.

SRE revisa semanalmente durante implantación y mensualmente en operación CPU, memoria, latencia, IOPS, edad y drenaje de colas y atraso de réplica. Arquitectura incorpora cambios al repositorio, T-11 y al inventario económico. La prueba cubre normal, peak, crecimiento, una tarea y una zona retiradas, emergencia y reconexión. La reserva local del 40\,\% se mantiene; antes de habilitar crecimiento se redimensiona el nodo o se demuestra el mismo margen con la nueva mezcla, sin atribuir a ocho núcleos capacidad medida para cualquier carga. El CLIENTE recibe proyección y consumo por ambiente/unidad, con alertas de presupuesto al 80 y 100\,\% del presupuesto autorizado; el aumento previsto por encima del 20\,\% activa revisión. Precios, tarifas, montos, comparación valorizada Fargate/EC2 y evidencia de calculadora se conservan exclusivamente en la Oferta Económica, bajo los mismos identificadores, conforme a BA 50.2; una alerta de gasto no detiene la captura crítica \parencites[RT-02.10, p.~7; RT-03.06, p.~8; RT-15.01, p.~28]{pucv-transversal}[art.~50.2, p.~29]{pucv-admin}.
